\documentclass[11pt,a4paper]{article}
\usepackage[utf8]{inputenc}
\usepackage[vietnamese]{babel}
\usepackage[left=2cm,right=2cm,top=2cm,bottom=2cm]{geometry}
\usepackage{mathptmx}
\usepackage{enumitem}
\usepackage{tabularx}
\usepackage{booktabs}
\usepackage{array}

\pagestyle{plain}
\linespread{1.15}

\begin{document}


\begin{center}\textbf{\large CỘNG HÒA XÃ HỘI CHỦ NGHĨA VIỆT NAM}\end{center}
\vspace{0.2cm}
\textbf{Độc lập – Tự do – Hạnh phúc}\\[0.15cm]

\par

\vspace{0.2cm}\hrule\vspace{0.3cm}

\par
\textbf{CÔNG AN PHƯỜNG PHÂN KHU CẢNG}\\[0.15cm]
Số: \texttt{01/P-TNTB}\\[0.15cm]
\textit{Hà Nội, hồi 06 giờ 45 phút ngày 25 tháng 07 năm 2016}\\[0.15cm]

\par

\vspace{0.2cm}\hrule\vspace{0.3cm}

\par

\begin{center}\textbf{\large PHIẾU TIẾP NHẬN NGUỒN TIN VỀ TỘI PHẠM}\end{center}
\vspace{0.2cm}

\par

\vspace{0.2cm}\hrule\vspace{0.3cm}

\par

\subsection*{I. CÁN BỘ TIẾP NHẬN}
\begin{itemize}[leftmargin=1.2cm, itemsep=2pt]
  \item \textbf{Họ và tên cán bộ trực ban:} Thượng úy Đỗ Văn Tuấn.
  \item \textbf{Chức vụ:} Cán bộ Trực ban Hình sự — Công an Phường Phân khu Cảng.
\end{itemize}

\par

\vspace{0.2cm}\hrule\vspace{0.3cm}

\par

\subsection*{II. THÔNG TIN NGƯỜI BÁO TIN}
\begin{itemize}[leftmargin=1.2cm, itemsep=2pt]
  \item \textbf{Họ và tên:} \textbf{NGUYỄN THỊ LỤA} | \textbf{Năm sinh:} 1956.
  \item \textbf{Số CCCD:} \texttt{148923057}.
  \item \textbf{Nơi cư trú:} Số 12, Đường Bờ Sông, Phường Phân khu Cảng, TP. Hà Nội (nhà sát vách số 14).
  \item \textbf{Hình thức báo tin:} Trực tiếp đến trụ sở Công an Phường trình báo khẩn cấp lúc 06:45.
\end{itemize}

\par

\vspace{0.2cm}\hrule\vspace{0.3cm}

\par

\subsection*{III. TÓM TẮT NỘI DUNG TRÌNH BÁO BAN ĐẦU}
Khoảng 06 giờ 30 phút sáng nay (25/07/2016), bà Lụa ra quét ngõ thì nhìn thấy cổng sắt và cửa chính căn nhà số 14 Đường Bờ Sông (nhà anh Nguyễn Văn Khang) mở toang, đèn điện phòng khách bật sáng trưng.\\[0.15cm]

\par
Bà Lụa gọi nhưng không thấy ai trả lời. Khi nhìn vào giữa phòng khách, bà hoảng hốt phát hiện anh Nguyễn Văn Khang đang nằm bất động trên nền gạch, xung quanh loang một vũng máu sẫm màu rất lớn, bên cạnh có bộ bình trà bị đập vỡ vụn. Bà Lụa chạy lại lay gọi thì thấy người nạn nhân đã lạnh ngắt, toàn thân co cứng. Bà lập tức hô hoán hàng xóm xung quanh và chạy thẳng đến Công an Phường báo án.\\[0.15cm]

\par

\vspace{0.2cm}\hrule\vspace{0.3cm}

\par

\subsection*{IV. BIỆN PHÁP XỬ LÝ BAN ĐẦU}
\begin{enumerate}[label=\arabic*., leftmargin=1.2cm, itemsep=2pt]
  \item Báo cáo khẩn cấp Trưởng Công an Phường và Giám đốc Công an TP.
  \item Cử Tổ công tác gồm 03 đồng chí xuống bảo vệ, phong tỏa hiện trường số 14 Đường Bờ Sông, nghiêm cấm người dân hiếu kỳ ra vào làm xáo trộn dấu vết.
  \item Chuyển hồ sơ ban đầu lên Phòng Cảnh sát Hình sự (PC02) Công an TP. Hà Nội thụ lý theo thẩm quyền.
\end{enumerate}

\par

\vspace{0.2cm}\hrule\vspace{0.3cm}

\par
\textbf{NGƯỜI BÁO TIN}\\[0.15cm]
\textit{(Ký, ghi rõ họ tên)}\\[0.15cm]
\textbf{Nguyễn Thị Lụa}\\[0.15cm]

\par
\textbf{CÁN BỘ TIẾP NHẬN TIN BÁO}\\[0.15cm]
\textit{(Ký, ghi rõ họ tên)}\\[0.15cm]
\textbf{Thượng úy Đỗ Văn Tuấn}\\[0.15cm]

\par


\end{document}
